\section{Discussion}
\label{sec:analysis}

\subsection{Context coherence as the mediating variable}

The results above suggest a coherent explanatory framework: the missing
variable between retrieval quality and answer faithfulness is \emph{context
coherence}. Language models are pattern-completion systems---when the context
presents a coherent narrative with sufficient evidential overlap, the model
continues that narrative faithfully. When the context has gaps from
fragmentation, domain mismatch, or topic drift, the model fills those gaps
from parametric memory, producing hallucination.

Table~\ref{tab:coherence_corr} reports Spearman correlations between
coherence-oriented retrieval metrics and answer faithfulness.

\begin{table}[!htbp]
\centering
\caption{Spearman correlations between retrieval metrics and faithfulness.
Llama-3-8B, chunk = 1024, $n = 30$ per condition. Per-query correlations
are modest due to small samples. The stronger signal is at the condition
level: high coherence coincides with low hallucination.}
\label{tab:coherence_corr}
\footnotesize
\begin{tabular}{llllc}
\toprule
Dataset & Condition & Metric & $\rho$ & $p$ \\
\midrule
PubMedQA & Baseline & retrieval\_entropy  & 0.332 & 0.073 \\
PubMedQA & Baseline & mean\_jaccard       & 0.279 & 0.136 \\
PubMedQA & Baseline & sim\_spread         & 0.244 & 0.195 \\
\midrule
PubMedQA & HCPC-v2  & embedding\_variance & 0.336 & 0.069 \\
PubMedQA & HCPC-v2  & retrieval\_entropy  & 0.315 & 0.090 \\
PubMedQA & HCPC-v2  & sim\_spread         & 0.305 & 0.102 \\
\midrule
SQuAD    & HCPC-v2  & mean\_query\_chunk\_sim & 0.323 & 0.082 \\
SQuAD    & HCPC-v2  & ccs                     & 0.313 & 0.092 \\
SQuAD    & Baseline & mean\_query\_chunk\_sim & 0.210 & 0.265 \\
\bottomrule
\end{tabular}
\end{table}

The per-query correlations ($\rho \leq 0.34$, all $p > 0.05$) are modest
and individually non-significant, which is expected at $n = 30$ and
reflects a genuine limitation: at the individual query level, faithfulness
depends on many factors beyond context coherence (question difficulty,
answer extractability, model confidence calibration). CCS is not designed
as a per-query oracle. The consistent \emph{direction} of the correlations
across all nine condition-metric pairs (all positive for coherence
indicators, all in the predicted direction) is more informative than any
single $p$-value, suggesting a real but modest per-query signal that would
require larger samples to confirm statistically.

The stronger and more actionable signal is at the \emph{condition level}:
CCS = 1.0 on SQuAD coincides with 0\% hallucination, while CCS = 0.57 on
PubMedQA coincides with 20\% hallucination. We do not claim a validated
functional relationship from two condition-level observations; rather, these
endpoints are consistent with the per-query trends and with the mechanistic
explanation that coherent contexts reduce the generator's need to bridge
evidential gaps. Figure~\ref{fig:ccs_hallucination} captures this aggregate
relationship, and Figure~\ref{fig:proxy} (discussed below) shows that
CCS-based proxy scores align with full answer-level evaluation.

This interpretation explains two observations that would otherwise seem
unrelated. First, HCPC-v2 helps on SQuAD by preserving an already coherent
retrieval set (intervening rarely), whereas it helps only modestly on
PubMedQA because the weak coherence originates upstream in the encoder. No
amount of selective refinement can fully recover a corpus representation that
is misaligned with the domain. Second, chunk size dominates the between-factor variance
(90.1\% of variance among the three tested parameters) because larger
chunks mechanically preserve intra-document coherence: topic sentences,
supporting details, and transitions remain together instead of being split
across fragments.

\input{../figures/ccs_and_proxy}

\subsection{Proxy-based threshold tuning}

Phase~5 explored 168 threshold configurations per dataset. Running full
generation and NLI scoring for each would be prohibitively expensive.
Instead, we used a proxy score combining CCS, embedding variance, retrieval
entropy, and mean query-chunk similarity (weights: 0.4, $-0.3$, $-0.2$, 0.1
respectively). The weights were selected based on the direction and
approximate magnitude of each metric's correlation with faithfulness in
Phases~1--2, not through cross-validation; they should be treated as a
reasonable starting point rather than an optimized configuration. As shown
in Figure~\ref{fig:proxy}, this proxy identifies the same high-performing
operating regions as full answer-level evaluation: on SQuAD, the top-10
configurations by proxy score all achieve CCS = 1.0, matching the
zero-hallucination region identified by full evaluation
(Appendix~\ref{app:sweep}).

The practical value is that a deployment team can prune most of the threshold
search space using retrieval-time statistics---which require no
generation---and reserve expensive answer-level evaluation for a small
shortlist of configurations. This is especially valuable for iterative RAG
development, where the cost of evaluating every retrieval tweak through
end-to-end generation can dominate experimentation time.

\subsection{Model-specific sensitivity}

The variance decomposition reveals a finding with important practical
implications: different model architectures respond to retrieval
configuration parameters in fundamentally different ways. On PubMedQA,
Mistral's faithfulness variance is governed by top-$k$ (97.9\%), while
Llama-3's is governed by chunk size (95.2\%). This means that a single set
of ``RAG best practices'' does not transfer across model families. The
optimal chunking strategy for one generator may be suboptimal for another,
and configurations should be tuned per-model rather than applied universally.

\subsection{Deployment implications}

Three operational lessons emerge from the study:

\begin{enumerate}
  \item \textbf{Evaluate retrieval and generation quality separately.}
        Passage-level relevance is necessary but not sufficient. A higher
        retrieval score can yield a worse answer context when coherence is
        degraded. Deployment monitoring should track CCS or a similar
        coherence metric alongside retrieval precision.

  \item \textbf{Prioritize chunking over prompting.} Among the
        configuration parameters we tested, chunk size accounts for
        450$\times$ more faithfulness variance than prompt strategy on
        SQuAD. Engineering effort should focus on chunk size selection,
        overlap strategy, and document structure preservation before
        investing in prompt optimization.

  \item \textbf{Validate coherence before deploying in specialized domains.}
        If CCS is high, the retrieval set is likely safe and refinement is
        unnecessary. If CCS is low, determine whether the cause is
        individual weak passages (refinement may help) or systemic
        encoder-corpus misalignment (fix the embedding model or use
        domain-specific retrieval instead).
\end{enumerate}

\subsection{Strengths of this study}

Several design choices strengthen the internal validity of our findings.
The shared pipeline architecture eliminates confounding from deployment
differences. The phased evaluation design isolates individual variables
while enabling cross-phase comparisons. The inclusion of both a
domain-matched (SQuAD) and domain-mismatched (PubMedQA) setting tests
whether findings depend on retrieval quality. The threshold sweep across 168
configurations verifies robustness rather than reporting results at a single
operating point. The zero-shot baseline provides an anchor for assessing
whether RAG itself is beneficial.

\subsection{Limitations}
\label{sec:limitations}

We acknowledge several limitations that affect the scope of our claims:

\paragraph{Sample size.} The HCPC head-to-head comparison uses $n = 30$ per
condition (90 total queries). At this sample size, hallucination rates have
a resolution of 3.3~pp, meaning the 0\% vs.\ 10\% SQuAD difference
reflects exactly 0 vs.\ 3 hallucinated responses. While the effects we
report are large (Cohen's $d > 1.0$) and consistent with trends across
6,000+ queries in earlier phases, the per-condition sample limits the
precision of effect size estimates and prevents strong claims about small
effects. All per-query CCS correlations are individually non-significant
($p > 0.05$). Increasing Phase~6 to $n = 100$+ would tighten confidence
intervals on both the hallucination rates and the CCS-faithfulness
relationship.

\paragraph{Dataset coverage.} The headline two-dataset comparison covers
domain-matched (SQuAD) and domain-mismatched (PubMedQA) extremes.
Section~\ref{sec:multidataset} extends the evaluation to four additional
datasets (NaturalQuestions, TriviaQA, HotpotQA distractor split, and
FinanceBench) covering open-domain factoid retrieval, multi-hop
reasoning, and a specialized professional domain. Conversational QA and
open-ended generation remain outside scope.

\paragraph{Model scale.} The three generators evaluated in
§\ref{sec:multidataset} (Mistral-7B, Llama-3-8B, Qwen2.5-7B) all sit in
the 7--8B parameter range. Larger models may handle incoherent contexts
differently---for example, by being more robust to fragmentation or by
hallucinating more confidently. Our findings should be validated at
larger scales before being applied to production systems using 70B+
models.

\paragraph{External method comparison.} The original draft of this paper
omitted comparison against Self-RAG and CRAG. Section~\ref{sec:headtohead}
addresses that gap with a five-condition head-to-head on the same
multi-dataset query set, including a CRAG reimplementation (the published
T5 evaluator is private; we substitute a cross-encoder threshold and
document the substitution transparently) and the published Self-RAG-7B
checkpoint. The remaining limitation is that Self-RAG is trained on a
specific reflection-token vocabulary; an apples-to-apples comparison
against a Self-RAG variant tuned on our exact corpus is out of scope here
and noted as future work.

\paragraph{Embedding model.} All experiments use a single general-domain
encoder. The CCS values and refinement thresholds reported here may not
transfer directly to systems using domain-specific or larger embedding
models.

\paragraph{Faithfulness evaluation.} NLI-based scoring captures
entailment but may miss subtle hallucinations (e.g., correct facts not
supported by the context) or penalize faithful paraphrases that are
lexically distant from the context. Section~\ref{sec:human} reports a
500-item Prolific study (1{,}000 annotation tasks) that calibrates the
NLI scorer against double-annotated human ratings and tests whether the
coherence paradox replicates under human evaluation. The reporting
commitments in §\ref{sec:human} include the falsification scenario where
NLI and humans diverge.

\paragraph{CCS granularity.} CCS operates at the condition level, not as a
per-query oracle. The modest per-query correlations ($\rho \leq 0.34$)
reflect both sample size limitations and the genuine complexity of
faithfulness, which depends on factors beyond context coherence alone.

\subsection{When HCPC-v2 does not help}

Three specific failure scenarios emerged from our analysis:

\begin{enumerate}
  \item \textbf{Systemic encoder-corpus misalignment.} When the embedding
        model consistently produces low-coherence retrieval (as on PubMedQA),
        selective refinement can mitigate but not solve the problem. The
        correct intervention is a domain-specific encoder, not a better
        refinement policy.

  \item \textbf{Strong baseline retrieval.} When the baseline retrieval is
        already coherent (CCS $\approx$ 1.0), HCPC-v2 adds computational
        overhead without meaningful improvement. The method's value is
        precisely in settings where some queries produce weak retrieval sets
        while others do not.

  \item \textbf{Rank protection masking genuine weakness.} In 3 of 3
        hallucinated responses under HCPC-v2 in extended experiments, the
        failure occurred because the rank protection gate prevented
        refinement of a passage that genuinely needed it (it ranked in the
        top 2 despite being inadequate). This suggests that rank-based
        protection, while generally beneficial, can occasionally backfire.
\end{enumerate}
